% ===================================================================
%  Семинар 3. Содержательная часть — единый источник для обоих файлов.
%
%  Темы взяты из «Семинар 3.docx»: сравнение регрессий (R^2_adj, AIC, BIC),
%  F-тест, тест на линейные ограничения, тест Рамсея, фиктивные переменные,
%  мультиколлинеарность и VIF. Теоретические вопросы — «Тема 2_Задания
%  к семинарам.docx» (множественная регрессия в матричной форме).
%
%  Задачи 4.4, 4.6, 4.14, 4.15 — глава 4 «Различные аспекты множественной
%  регрессии»: условия по учебнику Магнуса–Катышева–Пересецкого, разборы —
%  по «Сборнику задач к начальному курсу эконометрики» (РАНХиГС), с. 107–126.
%
%  Раздел «Основные формулы» и всё внутри \begin{solution}...\end{solution}
%  попадает только в seminar_03_solutions.pdf. Блоки \begin{excel}...\end{excel}
%  печатаются в обоих файлах: это часть задания.
%
%  Числовые результаты задач 4.14 и 4.15 пересчитаны независимо и совпадают
%  с таблицами 4.7–4.9 и 4.11–4.13 сборника.
% ===================================================================

\seminartitle{Семинар 3}{Множественная регрессия: тесты, фиктивные переменные,\\ мультиколлинеарность}

% ===================================================================
% Теория печатается только в файле с решениями.
\ifsolutions
\section{Основные формулы}

% -------------------------------------------------------------------
\subsection{Модель множественной регрессии}\label{sec:model}

В матричной форме модель записывается как
\begin{equation}
  \mathbf{y} = \mathbf{X}\boldsymbol{\beta} + \boldsymbol{\varepsilon},
  \label{eq:model}
\end{equation}
где $\mathbf{y}$ — вектор наблюдений размера $n\times 1$, $\mathbf{X}$ — матрица
регрессоров размера $n\times k$ (первый столбец — единицы, если в модели есть
свободный член), $\boldsymbol{\beta}$ — вектор $k$ оцениваемых коэффициентов,
$\boldsymbol{\varepsilon}$ — вектор ошибок. Число $k$ считается \emph{вместе} со
свободным членом.

\textbf{Предпосылки классической модели:}
\begin{enumerate}[label=\arabic*), leftmargin=2.2em, itemsep=2pt]
  \item выполнена спецификация \eqref{eq:model};
  \item $\mathbf{X}$ детерминирована и имеет полный ранг $k$: столбцы линейно
        независимы и $n > k$;
  \item[3a)] $\E\boldsymbol{\varepsilon} = \mathbf{0}$;
  \item[3b)] $\D(\boldsymbol{\varepsilon}) = \sigma^2\mathbf{I}_n$ — ошибки
        гомоскедастичны и некоррелированы;
  \item[3c)] $\boldsymbol{\varepsilon} \sim N(\mathbf{0}, \sigma^2\mathbf{I}_n)$;
        нужна только для распределений тестовых статистик.
\end{enumerate}

\textbf{Геометрический вывод МНК-оценки.} Прогноз $\mathbf{X}\mathbf{b}$ при
всевозможных $\mathbf{b} \in \mathbb{R}^k$ пробегает линейную оболочку столбцов
матрицы $\mathbf{X}$ — подпространство $L(\mathbf{X}) \subset \mathbb{R}^n$
размерности $k$ (столбцы линейно независимы по предпосылке~2). МНК минимизирует
\[
  S(\mathbf{b}) = (\mathbf{y} - \mathbf{X}\mathbf{b})^{\top}(\mathbf{y} - \mathbf{X}\mathbf{b})
                = \bigl\|\mathbf{y} - \mathbf{X}\mathbf{b}\bigr\|^2 ,
\]
то есть \emph{квадрат расстояния} от точки $\mathbf{y}$ до этого подпространства.
Ближайшая к $\mathbf{y}$ точка подпространства — ортогональная проекция
$\mathbf{y}$ на него, а это значит, что вектор остатков должен быть
перпендикулярен подпространству, то есть каждому столбцу $\mathbf{X}$:
\begin{equation}
  \mathbf{X}^{\top}\mathbf{e} = \mathbf{0},
  \qquad
  \mathbf{e} = \mathbf{y} - \mathbf{X}\widehat{\boldsymbol{\beta}} .
  \label{eq:normal}
\end{equation}

\begin{center}
\begin{tikzpicture}[>=stealth, line join=round, scale=1.0]
  % подпространство L(X)
  \fill[black!5]   (-0.7,-0.55) -- (6.3,-0.55) -- (7.7,1.45) -- (0.7,1.45) -- cycle;
  \draw[black!45]  (-0.7,-0.55) -- (6.3,-0.55) -- (7.7,1.45) -- (0.7,1.45) -- cycle;
  \node[black!60, anchor=west] at (6.75,1.05) {$L(\mathbf{X})$};

  \coordinate (O) at (1.05,0.10);
  \coordinate (P) at (4.95,0.85);   % ровно под y: тогда e рисуется вертикально
  \coordinate (Y) at (4.95,3.35);
  \coordinate (B) at (6.15,0.40);

  % катет в плоскости и гипотенуза треугольника (y, y-с-крышкой, Xb)
  \draw[black!55]         (P) -- (B);
  \draw[dotted, black!55] (B) -- (Y);

  % прямой угол между e и плоскостью
  \draw[black!55]
    ($(P)!0.32cm!(B)$) -- ($($(P)!0.32cm!(B)$)+($(P)!0.32cm!(Y)$)-(P)$)
    -- ($(P)!0.32cm!(Y)$);

  % векторы
  \draw[->, very thick]           (O) -- (Y) node[anchor=south west] {$\mathbf{y}$};
  \draw[->, very thick, black!70] (O) -- (P)
        node[anchor=north, pos=1, yshift=-3pt] {$\widehat{\mathbf{y}}$};
  \draw[->, very thick, dashed]   (P) -- (Y)
        node[anchor=west, midway] {$\mathbf{e}$};
  \draw[->, black!55]             (O) -- (B)
        node[anchor=north, yshift=-2pt] {$\mathbf{X}\mathbf{b}$};

  \fill (O) circle (1.4pt);
  \fill (P) circle (1.4pt);
  \fill (Y) circle (1.4pt);
  \fill[black!55] (B) circle (1.2pt);
\end{tikzpicture}

\smallskip
{\small\itshape Проекция $\widehat{\mathbf{y}}$ — ближайшая к $\mathbf{y}$ точка
подпространства $L(\mathbf{X})$; остаток $\mathbf{e}$ перпендикулярен ему.
Для любой другой точки $\mathbf{X}\mathbf{b}$ отрезок до $\mathbf{y}$ —
гипотенуза прямоугольного треугольника, поэтому он длиннее.}
\end{center}

\emph{Почему ортогональность даёт минимум.} Возьмём произвольное $\mathbf{b}$ и
обозначим $\mathbf{d} = \widehat{\boldsymbol{\beta}} - \mathbf{b}$. Тогда
\[
  \mathbf{y} - \mathbf{X}\mathbf{b}
  = \bigl(\mathbf{y} - \mathbf{X}\widehat{\boldsymbol{\beta}}\bigr)
    + \mathbf{X}\bigl(\widehat{\boldsymbol{\beta}} - \mathbf{b}\bigr)
  = \mathbf{e} + \mathbf{X}\mathbf{d},
\]
причём слагаемые ортогональны: $(\mathbf{X}\mathbf{d})^{\top}\mathbf{e}
= \mathbf{d}^{\top}(\mathbf{X}^{\top}\mathbf{e}) = 0$ по \eqref{eq:normal}. По теореме
Пифагора
\[
  \bigl\|\mathbf{y} - \mathbf{X}\mathbf{b}\bigr\|^2
  = \|\mathbf{e}\|^2 + \|\mathbf{X}\mathbf{d}\|^2 \ \ge\ \|\mathbf{e}\|^2 ,
\]
и равенство достигается только при $\mathbf{X}\mathbf{d} = \mathbf{0}$, то есть
(снова по предпосылке~2) при $\mathbf{d} = \mathbf{0}$. Значит, минимум
единственный и достигается в точке, задаваемой условием \eqref{eq:normal}.

Осталось раскрыть \eqref{eq:normal}. Подставив
$\mathbf{e} = \mathbf{y} - \mathbf{X}\widehat{\boldsymbol{\beta}}$, получаем
\emph{нормальные уравнения}:
\[
  \mathbf{X}^{\top}\bigl(\mathbf{y} - \mathbf{X}\widehat{\boldsymbol{\beta}}\bigr) = \mathbf{0}
  \qquad\Longleftrightarrow\qquad
  \mathbf{X}^{\top}\mathbf{X}\widehat{\boldsymbol{\beta}} = \mathbf{X}^{\top}\mathbf{y} .
\]
Матрица $\mathbf{X}^{\top}\mathbf{X}$ обратима, поскольку $\mathbf{X}$ имеет полный
ранг, — умножая на $(\mathbf{X}^{\top}\mathbf{X})^{-1}$, получаем \textbf{МНК-оценку
вектора коэффициентов}:
\begin{equation}
  \boxed{\ \widehat{\boldsymbol{\beta}} = (\mathbf{X}^{\top}\mathbf{X})^{-1}\mathbf{X}^{\top}\mathbf{y}\ }
  \label{eq:ols}
\end{equation}
Никакого дифференцирования по вектору здесь не понадобилось: всё свелось к тому,
что кратчайшее расстояние до подпространства даёт перпендикуляр.

\textbf{Два следствия.} Во-первых, если первый столбец $\mathbf{X}$ состоит из
единиц, то одно из равенств \eqref{eq:normal} — это в точности
$\sum e_t = 0$: остатки в сумме дают нуль. Во-вторых, разложение
$\TSS = \RSS + \ESS$ — та же теорема Пифагора, только применённая к отклонениям
от среднего: вектор с координатами $Y_t - \overline{Y}$ раскладывается на
подогнанную часть с координатами $\widehat{Y}_t - \overline{Y}$ и
перпендикулярный ей остаток $\mathbf{e}$.

\textbf{Математическое ожидание и дисперсия оценки.} Подставив \eqref{eq:model}
в \eqref{eq:ols}, выразим оценку через истинные коэффициенты и ошибки модели:
\begin{equation}
  \widehat{\boldsymbol{\beta}}
  = (\mathbf{X}^{\top}\mathbf{X})^{-1}\mathbf{X}^{\top}
    \bigl(\mathbf{X}\boldsymbol{\beta} + \boldsymbol{\varepsilon}\bigr)
  = \boldsymbol{\beta} + \mathbf{A}\boldsymbol{\varepsilon},
  \qquad
  \mathbf{A} = (\mathbf{X}^{\top}\mathbf{X})^{-1}\mathbf{X}^{\top} .
  \label{eq:beta-err}
\end{equation}
Матрица $\mathbf{A}$ детерминирована (предпосылка~2), поэтому её выносят за знак
матожидания как постоянный множитель. Отсюда сразу несмещённость:
\[
  \E\widehat{\boldsymbol{\beta}}
  = \boldsymbol{\beta} + \mathbf{A}\,\E\boldsymbol{\varepsilon}
  = \boldsymbol{\beta}
  \qquad\text{(предпосылка 3a)} .
\]
Для вектора дисперсия — это матрица ковариаций: на диагонали стоят дисперсии
координат, вне диагонали — их попарные ковариации,
\[
  \D(\mathbf{z}) = \E\bigl[(\mathbf{z} - \E\mathbf{z})(\mathbf{z} - \E\mathbf{z})^{\top}\bigr],
  \qquad\text{в частности}\qquad
  \E\bigl[\boldsymbol{\varepsilon}\boldsymbol{\varepsilon}^{\top}\bigr]
  = \D(\boldsymbol{\varepsilon}) = \sigma^2\mathbf{I}_n
\]
по предпосылкам 3a и 3b. Подставляя \eqref{eq:beta-err} и учитывая, что при
транспонировании произведения порядок сомножителей меняется,
$(\mathbf{A}\boldsymbol{\varepsilon})^{\top} = \boldsymbol{\varepsilon}^{\top}\mathbf{A}^{\top}$,
а $(\mathbf{X}^{\top}\mathbf{X})^{-1}$ симметрична, а значит
$\mathbf{A}^{\top} = \mathbf{X}(\mathbf{X}^{\top}\mathbf{X})^{-1}$:
\[
  \begin{aligned}
    \D(\widehat{\boldsymbol{\beta}})
    &= \E\bigl[(\widehat{\boldsymbol{\beta}} - \E\widehat{\boldsymbol{\beta}})
               (\widehat{\boldsymbol{\beta}} - \E\widehat{\boldsymbol{\beta}})^{\top}\bigr]
     = \E\bigl[(\widehat{\boldsymbol{\beta}} - \boldsymbol{\beta})
               (\widehat{\boldsymbol{\beta}} - \boldsymbol{\beta})^{\top}\bigr] \\[4pt]
    &= \E\bigl[(\mathbf{A}\boldsymbol{\varepsilon})(\mathbf{A}\boldsymbol{\varepsilon})^{\top}\bigr]
     = \E\bigl[\mathbf{A}\boldsymbol{\varepsilon}\boldsymbol{\varepsilon}^{\top}\mathbf{A}^{\top}\bigr]
     = \mathbf{A}\,\E\bigl[\boldsymbol{\varepsilon}\boldsymbol{\varepsilon}^{\top}\bigr]\mathbf{A}^{\top}
     = \sigma^2\mathbf{A}\mathbf{A}^{\top} \\[4pt]
    &= \sigma^2 (\mathbf{X}^{\top}\mathbf{X})^{-1}\mathbf{X}^{\top}
       \mathbf{X}(\mathbf{X}^{\top}\mathbf{X})^{-1}
     = \sigma^2 (\mathbf{X}^{\top}\mathbf{X})^{-1} .
  \end{aligned}
\]
Итого при предпосылках 1--3b
\begin{equation}
  \begin{gathered}
    \E\widehat{\boldsymbol{\beta}} = \boldsymbol{\beta},
    \qquad
    \D(\widehat{\boldsymbol{\beta}}) = \sigma^2(\mathbf{X}^{\top}\mathbf{X})^{-1},
    \\[6pt]
    s^2 = \frac{\ESS}{n-k}
        = \frac{\sum e_t^2}{n-k}
        = \frac{\mathbf{e}^{\top}\mathbf{e}}{n-k},
    \qquad
    \E s^2 = \sigma^2 .
  \end{gathered}
  \label{eq:props}
\end{equation}
На диагонали матрицы $\D(\widehat{\boldsymbol{\beta}})$ стоят дисперсии отдельных
коэффициентов:
\[
  \D(\widehat{\beta}_j)
  = \sigma^2\bigl[(\mathbf{X}^{\top}\mathbf{X})^{-1}\bigr]_{jj} ,
\]
так что, заменив неизвестную $\sigma^2$ её несмещённой оценкой $s^2$, получаем
стандартные ошибки коэффициентов.
По теореме Гаусса--Маркова МНК-оценка эффективна в классе линейных несмещённых
оценок (семинар~2 — парный случай). Стандартная ошибка $j$-го коэффициента и
$t$-статистика:
\begin{equation}
  s_{\widehat{\beta}_j} = s\sqrt{\bigl[(\mathbf{X}^{\top}\mathbf{X})^{-1}\bigr]_{jj}},
  \qquad
  t = \frac{\widehat{\beta}_j - \beta_j^0}{s_{\widehat{\beta}_j}} \sim t(n-k)
  \quad\text{при } H_0\colon \beta_j = \beta_j^0 .
  \label{eq:t}
\end{equation}
Степеней свободы $n-k$: из $n$ наблюдений вычитается $k$ оценённых коэффициентов.

% -------------------------------------------------------------------
\subsection{Качество подгонки и сравнение моделей}\label{sec:fit}

При наличии свободного члена сохраняется разложение сумм квадратов
\begin{equation}
  \underbrace{\sum (Y_t - \overline{Y})^2}_{\TSS}
  = \underbrace{\sum (\widehat{Y}_t - \overline{Y})^2}_{\RSS}
  + \underbrace{\sum (Y_t - \widehat{Y}_t)^2}_{\ESS},
  \qquad \ESS = \sum e_t^2 = \mathbf{e}^{\top}\mathbf{e},
  \label{eq:ss}
\end{equation}
и на нём построен коэффициент детерминации
\begin{equation}
  R^2 = \frac{\RSS}{\TSS} = 1 - \frac{\ESS}{\TSS} .
  \label{eq:r2}
\end{equation}
Здесь $\TSS$ зависит только от $\mathbf{y}$, то есть от данных, а не от модели, —
при добавлении регрессора он не меняется. А $\ESS$ — это минимум квадрата
расстояния от $\mathbf{y}$ до $L(\mathbf{X})$, и добавление столбца
$\mathbf{x}_{k+1}$ расширяет подпространство, так что минимум берётся по более
широкому множеству:
\[
  \ESS_{k+1}
  = \min_{\mathbf{b},\,c} \bigl\|\mathbf{y} - \mathbf{X}\mathbf{b} - c\,\mathbf{x}_{k+1}\bigr\|^2
  \ \le\ \bigl\|\mathbf{y} - \mathbf{X}\widehat{\boldsymbol{\beta}} - 0\cdot\mathbf{x}_{k+1}\bigr\|^2
  = \bigl\|\mathbf{y} - \mathbf{X}\widehat{\boldsymbol{\beta}}\bigr\|^2
  = \ESS_k :
\]
старая подгонка по-прежнему доступна — она получается при $c = 0$, — поэтому
новый минимум не может быть больше. Отсюда
\[
  R^2_{k+1} = 1 - \frac{\ESS_{k+1}}{\TSS} \ \ge\ 1 - \frac{\ESS_k}{\TSS} = R^2_k ,
\]
причём равенство лишь в вырожденном случае $\widehat{c} = 0$. Значит, $R^2$ не
убывает от добавления \emph{любого} регрессора, даже бессмысленного, и сравнивать
по нему модели с разным числом регрессоров нельзя. Для сравнения используют
скорректированный коэффициент детерминации и информационные критерии:
\begin{gather}
  R^2_{\text{adj}} = 1 - \frac{\ESS/(n-k)}{\TSS/(n-1)}
                   = 1 - \frac{n-1}{n-k}\bigl(1 - R^2\bigr)
                   = R^2 - \frac{k-1}{n-k}\bigl(1 - R^2\bigr),
  \label{eq:r2adj} \\[2pt]
  \mathrm{AIC} = \frac{2k}{n} + \ln\frac{\ESS}{n},
  \qquad
  \mathrm{BIC} = \frac{k\ln n}{n} + \ln\frac{\ESS}{n} .
  \label{eq:aic}
\end{gather}
В \eqref{eq:r2adj} каждая сумма квадратов делится на свои степени свободы, так
что добавление регрессора увеличивает $R^2_{\text{adj}}$ только тогда, когда
выигрыш в $\ESS$ перевешивает потерю степени свободы. Критерии AIC и BIC
устроены так же: первое слагаемое — штраф за число параметров, второе —
качество подгонки; лучше та модель, у которой критерий \emph{меньше}. У BIC
штраф строже, поэтому он выбирает более экономные модели.

% -------------------------------------------------------------------
\subsection{Проверка гипотез: $t$- и $F$-тесты}\label{sec:tests}

Отдельный коэффициент проверяется $t$-статистикой \eqref{eq:t}: гипотеза
$H_0\colon \beta_j = 0$ отвергается на уровне 5\,\%, если
$|t| > t_{0{,}025}(n-k)$.

Значимость регрессии \emph{в целом} — это гипотеза
$H_0\colon \beta_2 = \dots = \beta_k = 0$ (все коэффициенты, кроме свободного
члена, равны нулю). Её проверяют статистикой
\begin{equation}
  F = \frac{\RSS/(k-1)}{\ESS/(n-k)}
    = \frac{R^2}{1-R^2}\cdot\frac{n-k}{k-1} \sim F(k-1,\ n-k),
  \label{eq:f}
\end{equation}
второе равенство получается делением числителя и знаменателя на $\TSS$.
Гипотеза отвергается при $F > F_{0{,}05}(k-1, n-k)$. В парной регрессии
($k = 2$) $F = t^2$, и оба теста совпадают.

\textbf{$P$-значение.} Это вероятность получить статистику не менее
экстремальную, чем наблюдённая, \emph{при верной} $H_0$:
\[
  P = \Pr\bigl\{|t(n-k)| > |t_{\text{набл}}|\bigr\},
  \qquad
  P = \Pr\bigl\{F(k-1,\, n-k) > F_{\text{набл}}\bigr\}
\]
для $t$- и $F$-теста соответственно ($t$-тест двусторонний, поэтому берётся
модуль; для $F$-теста против $H_0$ говорят только большие значения). Правило
решения:
\[
  P < 0{,}05 \ \Longrightarrow\ H_0 \text{ отвергается на } 5\,\%\text{-ном уровне},
  \qquad
  P \geqslant 0{,}05 \ \Longrightarrow\ \text{не отвергается},
\]
и так же для любого другого уровня значимости: гипотеза отвергается, если
$P$ меньше него. Это то же самое сравнение, что и с критическим значением,
только записанное с другой стороны: вместо «на сколько статистика превысила
порог» смотрят, «насколько редким был бы такой результат при $H_0$». Удобство в
том, что одно число сразу даёт ответ для всех уровней: при $P = 0{,}03$ гипотеза
отвергается на 5\,\%-ном уровне и не отвергается на 1\,\%-ном. Именно
$P$-значения печатают эконометрические пакеты — столбец \emph{Probability} в
таблицах решений ниже.

\begin{remark}
$P$-значение — это \emph{не} вероятность того, что $H_0$ верна, и не мера силы
связи. Малое $P$ означает лишь, что при верной $H_0$ такие данные встречались бы
редко. Коэффициент может быть значим ($P$ мало) и при этом ничтожно мал по
величине: статистическая значимость и содержательная важность — разные вещи.
\end{remark}

% -------------------------------------------------------------------
\subsection{Тест на линейные ограничения}\label{sec:restr}

Пусть на коэффициенты наложено $q$ линейных ограничений (например, часть
коэффициентов равна нулю). Оценим модель дважды: без ограничений (UR) и с
ограничениями (R). Так как ограничение сужает множество допустимых
коэффициентов, всегда $\ESS_{\mathrm{R}} \ge \ESS_{\mathrm{UR}}$. Вопрос в том,
велик ли прирост:
\begin{equation}
  F = \frac{\bigl(\ESS_{\mathrm{R}} - \ESS_{\mathrm{UR}}\bigr)/q}
           {\ESS_{\mathrm{UR}}/(n-k)} \sim F(q,\ n-k),
  \label{eq:restr}
\end{equation}
где $k$ — число параметров модели \emph{без} ограничений. Если обе модели
оценены по одним и тем же данным, $\TSS$ у них одинакова, и \eqref{eq:restr}
можно записать через коэффициенты детерминации:
\[
  F = \frac{\bigl(R^2_{\mathrm{UR}} - R^2_{\mathrm{R}}\bigr)/q}
           {\bigl(1 - R^2_{\mathrm{UR}}\bigr)/(n-k)} .
\]
Частные случаи: при $q = k-1$ (все коэффициенты, кроме константы, равны нулю)
формула \eqref{eq:restr} превращается в \eqref{eq:f}; при $q = 1$ получается
$F = t^2$.

% -------------------------------------------------------------------
\subsection{Тест Чоу на структурный сдвиг}\label{sec:chow}

Пусть выборка разбита на две части объёмом $n_1$ и $n_2$ ($n_1 + n_2 = n$), и
проверяется гипотеза $H_0$ о том, что коэффициенты в обеих частях одинаковы.
Оценим регрессию трижды: по всей выборке ($\ESS_{\mathrm{R}}$) и отдельно по
каждой части ($\ESS_1$ и $\ESS_2$). Тогда
\begin{equation}
  F = \frac{\bigl(\ESS_{\mathrm{R}} - (\ESS_1 + \ESS_2)\bigr)/k}
           {(\ESS_1 + \ESS_2)/(n - 2k)} \sim F(k,\ n-2k).
  \label{eq:chow}
\end{equation}

\begin{remark}
Тест Чоу — частный случай \eqref{eq:restr}. Введём фиктивную переменную $D$,
равную единице для второй подвыборки, и оценим модель со всеми регрессорами,
умноженными на $D$: $\mathbf{y} = \mathbf{X}\boldsymbol{\beta} +
(D\cdot\mathbf{X})\boldsymbol{\gamma} + \boldsymbol{\varepsilon}$. Такая модель
подгоняет обе подвыборки независимо, поэтому её $\ESS$ равна $\ESS_1 + \ESS_2$,
а гипотеза $H_0\colon \boldsymbol{\gamma} = \mathbf{0}$ — это $q = k$
ограничений. Подстановка в \eqref{eq:restr} даёт в точности \eqref{eq:chow}.
\end{remark}

% -------------------------------------------------------------------
\subsection{Фиктивные переменные}\label{sec:dummy}

Фиктивная (дамми) переменная $d$ принимает значения 0 и 1 и позволяет включить
в регрессию качественный признак:
\begin{equation}
  y_t = \alpha + \beta x_t + \gamma d_t + \varepsilon_t
  \qquad\text{и}\qquad
  y_t = \alpha + \beta x_t + \gamma d_t + \delta\,(d_tx_t) + \varepsilon_t .
  \label{eq:dummy}
\end{equation}
В первой модели признак сдвигает \emph{свободный член}: при $d = 1$ прямая
поднимается на $\gamma$, наклон общий. Во второй добавлено \emph{взаимодействие}
$d\cdot x$, поэтому меняется и наклон: он равен $\beta$ при $d = 0$ и
$\beta + \delta$ при $d = 1$. Гипотеза «признак ничего не меняет» — это
$H_0\colon \gamma = \delta = 0$, то есть $q = 2$ ограничения в \eqref{eq:restr}.

\emph{Ловушка фиктивных переменных.} Если признак принимает $m$ значений, в
модель со свободным членом включают только $m-1$ фиктивную переменную. Иначе их
сумма равна столбцу единиц, столбцы $\mathbf{X}$ линейно зависимы, предпосылка~2
нарушена и оценка \eqref{eq:ols} не существует. Пропущенная категория
становится \emph{базовой}, а коэффициенты показывают отклонение от неё.

\emph{Сезонность.} Для квартальных данных вводят три фиктивные переменные
$\mathit{QRT}1$, $\mathit{QRT}2$, $\mathit{QRT}3$ (четвёртый квартал — база).
Проверка наличия сезонности — это $F$-тест \eqref{eq:restr} совместной
значимости всех трёх коэффициентов, а не три отдельных $t$-теста.

\emph{Перелом наклона в точке.} Если наклон меняется при переходе через
пороговое значение $x^{*}$, берут переменную $(x_t - x^{*})s_t$, где $s_t = 1$
при $x_t > x^{*}$ и $s_t = 0$ иначе. Тогда наклон равен $\beta$ слева от порога
и $\beta + \gamma$ справа, а сама функция в точке $x^{*}$ не разрывается.

% -------------------------------------------------------------------
\subsection{Тест Рамсея (RESET)}\label{sec:reset}

Тест проверяет, не потеряна ли нелинейность, то есть верна ли функциональная
форма. Оценивают исходную регрессию, сохраняют подогнанные значения
$\widehat{y}_t$ и оценивают расширенную модель
\begin{equation}
  y_t = \mathbf{x}_t^{\top}\boldsymbol{\beta}
        + \gamma_1\widehat{y}_t^{\,2} + \dots + \gamma_p\widehat{y}_t^{\,p+1}
        + \varepsilon_t ,
  \label{eq:reset}
\end{equation}
после чего проверяют $H_0\colon \gamma_1 = \dots = \gamma_p = 0$ по
\eqref{eq:restr} со статистикой $F(p,\ n-k-p)$. Отвержение $H_0$ означает, что
линейная спецификация неадекватна: степени $\widehat{y}_t$ подхватывают
пропущенную нелинейность.

% -------------------------------------------------------------------
\subsection{Мультиколлинеарность и VIF}\label{sec:vif}

\emph{Точная} мультиколлинеарность — линейная зависимость столбцов
$\mathbf{X}$ — нарушает предпосылку~2: матрица $\mathbf{X}^{\top}\mathbf{X}$
вырождена и оценка \eqref{eq:ols} не существует. На практике опаснее
\emph{сильная, но не точная} связь регрессоров: оценки существуют, но
\begin{itemize}[leftmargin=2.2em, itemsep=2pt]
  \item дисперсии \eqref{eq:props} велики, оценки неустойчивы — небольшое
        изменение данных сильно меняет коэффициенты;
  \item отдельные $t$-статистики малы, хотя $F$-тест \eqref{eq:f} уверенно
        отвергает незначимость регрессии в целом;
  \item знаки коэффициентов могут противоречить содержательному смыслу.
\end{itemize}
Диагностика — \emph{фактор инфляции дисперсии}. Для $j$-го регрессора оценивают
вспомогательную регрессию этого регрессора на все остальные и берут её
коэффициент детерминации $R_j^2$:
\begin{equation}
  \mathrm{VIF}_j = \frac{1}{1 - R_j^2},
  \qquad
  \D(\widehat{\beta}_j)
  = \mathrm{VIF}_j \cdot \frac{\sigma^2}{\sum_t (x_{tj} - \overline{x}_j)^2} .
  \label{eq:vif}
\end{equation}
Вторая формула объясняет название: $\mathrm{VIF}_j$ показывает, во сколько раз
дисперсия оценки больше, чем была бы при регрессорах, не связанных друг с
другом. Обычное эмпирическое правило: $\mathrm{VIF}_j > 10$, то есть
$R_j^2 > 0{,}9$, — признак серьёзной мультиколлинеарности.
\fi

% ===================================================================
\section{Задачи}

% -------------------------------------------------------------------
\problem{4.4}

Предположим, что вы оцениваете линейную функцию потребления
$c_t = \alpha + \beta y_t + \varepsilon_t$ среди $n$ индивидуумов. Как учесть
возможный сдвиг этой функции при переходе от городского к сельскому
потребителю, если вы считаете, что предельная (маржинальная) склонность к
потреблению постоянна, в то время как средняя склонность к потреблению может
меняться? Как проверить гипотезу о том, что предельные склонности к потреблению
индивидуумов с доходом выше и ниже уровня $y^{*}$ отличаются?

\begin{solution}
\textbf{1) Сдвиг при переходе от города к селу.} Предельная склонность к
потреблению — это прирост потребления на единицу дохода, то есть
$\partial c/\partial y = \beta$. Средняя склонность — доля дохода, идущая на
потребление:
\[
  \frac{c}{y} = \frac{\alpha}{y} + \beta ,
\]
она зависит не только от $\beta$, но и от $\alpha$. Значит, требование «$\beta$
постоянна, а средняя склонность может меняться» означает, что при переходе от
города к селу меняется только свободный член. По \eqref{eq:dummy} вводим
фиктивную переменную $r_t = 0$ для городского потребителя и $r_t = 1$ для
сельского:
\[
  c_t = \alpha + \beta y_t + \gamma r_t + \varepsilon_t .
\]

\begin{center}
\begin{tikzpicture}[>=stealth, line join=round, x=0.9cm, y=0.9cm]
  % параметры картинки: alpha = 1, beta = 0.5, gamma = 1.2, y0 = 3
  \draw[->] (0,0) -- (8.6,0) node[anchor=north] {$y$};
  \draw[->] (0,0) -- (0,6.6) node[anchor=east]  {$c$};
  \node[anchor=north east] at (0,0) {$0$};

  % две функции потребления: город (r = 0) и село (r = 1)
  \draw[very thick, accent]  (0,1.0) -- (8,5.0)
        node[anchor=west] {город, $r_t = 0$};
  \draw[very thick, solbar]  (0,2.2) -- (8,6.2)
        node[anchor=west] {село, $r_t = 1$};

  % свободные члены и сдвиг gamma
  \fill[accent] (0,1.0) circle (1.5pt);
  \fill[solbar] (0,2.2) circle (1.5pt);
  \node[anchor=east, accent] at (0,1.0) {$\alpha$};
  \node[anchor=east, solbar] at (0,2.2) {$\alpha + \gamma$};
  \draw[<->, thick] (1.0,1.5) -- (1.0,2.7)
        node[midway, anchor=west] {$\gamma$};

  % одинаковый наклон: треугольники приращений (у села — над прямой,
  % у города — под ней, чтобы подписи не налезали на соседнюю линию)
  \draw[black!60] (4.5,4.45) -- (4.5,5.45) -- (6.5,5.45);
  \node[anchor=east,  black!70] at (4.5,4.95) {\small $\beta\,\Delta y$};
  \node[anchor=south, black!70] at (5.5,5.45) {\small $\Delta y$};
  \draw[black!60] (5.0,3.5) -- (7.0,3.5) -- (7.0,4.5);
  \node[anchor=north, black!70] at (6.0,3.5) {\small $\Delta y$};
  \node[anchor=west,  black!70] at (7.0,4.0) {\small $\beta\,\Delta y$};

  % средняя склонность при одном и том же доходе y0 — наклоны лучей из нуля
  \draw[dotted, black!55] (3,0) -- (3,3.7);
  \node[anchor=north] at (3,0) {$y_0$};
  \draw[dashed, accent]  (0,0) -- (3,2.5);
  \draw[dashed, solbar]  (0,0) -- (3,3.7);
  \fill[accent] (3,2.5) circle (1.5pt);
  \fill[solbar] (3,3.7) circle (1.5pt);
\end{tikzpicture}

\smallskip
{\small\itshape Сплошные линии — функции потребления. Наклон у них общий,
$\beta$: предельная склонность одинакова (треугольники приращений равны).
Сельская прямая сдвинута на $\gamma$. Пунктирные лучи из начала координат
к точкам при одном и том же доходе $y_0$ имеют наклоны $c/y$, то есть это
средние склонности к потреблению. Из-за сдвига они различаются:
$(\alpha + \gamma)/y_0 + \beta$ против $\alpha/y_0 + \beta$.\par}
\end{center}

Наличие сдвига — это гипотеза $H_0\colon \gamma = 0$, которая проверяется
$t$-статистикой \eqref{eq:t} с $n - 3$ степенями свободы.

\medskip
\textbf{2) Разные предельные склонности выше и ниже уровня $y^{*}$.} Теперь
меняться должен наклон, причём только для индивидуумов с доходом выше $y^{*}$.
По п.~\ref{sec:dummy} вводим $s_t = 0$ при $y_t \leqslant y^{*}$ и $s_t = 1$ при
$y_t > y^{*}$ и оцениваем
\[
  c_t = \alpha + \beta y_t + \gamma\,(y_t - y^{*})s_t + \varepsilon_t .
\]
Предельная склонность к потреблению равна
\[
  \frac{\partial c}{\partial y} =
  \begin{cases}
    \beta, & y \leqslant y^{*}, \\[2pt]
    \beta + \gamma, & y > y^{*},
  \end{cases}
\]
поэтому проверяемая гипотеза — снова $H_0\colon \gamma = 0$. Множитель
$(y_t - y^{*})$ вместо просто $y_t$ выбран для того, чтобы функция потребления
не имела разрыва в точке $y^{*}$: при $y_t = y^{*}$ добавка обращается в нуль.
\end{solution}

% -------------------------------------------------------------------
\problem{4.6}

На основе квартальных данных с 1971 по 1976 г. с помощью метода наименьших
квадратов получено следующее уравнение:
\[
  y_t = \underset{(2{,}14)}{1{,}12}
      - \underset{(0{,}0034)}{0{,}0098}\,x_{t1}
      - \underset{(3{,}42)}{5{,}62}\,x_{t2}
      + \underset{(0{,}009)}{0{,}044}\,x_{t3},
\]
в скобках указаны стандартные ошибки, $\RSS = 110{,}32$, $\ESS = 21{,}43$.
\begin{enumerate}[label=\asbuk*), leftmargin=2.2em, itemsep=3pt]
  \item Проверьте значимость каждого из коэффициентов.
  \item Найдите коэффициент детерминации.
  \item Протестируйте значимость регрессии в целом.
  \item Когда в уравнение были добавлены три фиктивные переменные,
        соответствующие трём первым кварталам года, величина $\RSS$ выросла
        до 118,20. Проверьте гипотезу о наличии сезонности, сформулировав
        необходимые предположения о виде этой сезонности.
  \item Для той же исходной модели были раздельно проведены две регрессии на
        основе данных: 1-й квартал 1971 г. — 1-й квартал 1975 г. и 2-й квартал
        1975 г. — 4-й квартал 1976 г. соответственно, и получены следующие
        значения сумм квадратов остатков: $\ESS_1 = 12{,}25$,
        $\ESS_2 = 2{,}32$. Проверьте гипотезу о том, что между 1-м и 2-м
        кварталами 1975 г. произошло структурное изменение.
\end{enumerate}

\begin{solution}
Здесь $n = 24$ (шесть лет по четыре квартала) и $k = 4$, поэтому $n - k = 20$.

\medskip
\textbf{а)} Для каждого коэффициента проверяем $H_0\colon \beta_i = 0$ по
\eqref{eq:t}, сравнивая $|t_i| = |\widehat{\beta}_i|/s_{\widehat{\beta}_i}$ с
$t_{0{,}025}(20) = 2{,}086$:
\[
  |t_1| = \frac{1{,}12}{2{,}14} = 0{,}523,
  \qquad
  |t_2| = \frac{0{,}0098}{0{,}0034} = 2{,}882,
\]
\[
  |t_3| = \frac{5{,}62}{3{,}42} = 1{,}643,
  \qquad
  |t_4| = \frac{0{,}044}{0{,}009} = 4{,}889 .
\]
Значимы второй и четвёртый коэффициенты (при $x_{t1}$ и $x_{t3}$); свободный
член и коэффициент при $x_{t2}$ на 5\,\%-ном уровне незначимы.

\medskip
\textbf{б)} По \eqref{eq:r2}, учитывая $\TSS = \RSS + \ESS = 110{,}32 + 21{,}43
= 131{,}75$:
\[
  R^2 = \frac{\RSS}{\TSS} = \frac{110{,}32}{131{,}75} = 0{,}837 .
\]

\medskip
\textbf{в)} Гипотеза $H_0\colon \beta_2 = \beta_3 = \beta_4 = 0$ проверяется по
\eqref{eq:f}:
\[
  F = \frac{\RSS}{\ESS}\cdot\frac{n-k}{k-1}
    = \frac{110{,}32}{21{,}43}\cdot\frac{20}{3}
    = 34{,}319 > 3{,}098 = F_{0{,}05}(3, 20),
\]
то есть регрессия в целом значима.

\medskip
\textbf{г)} Три фиктивные переменные для первых трёх кварталов (четвёртый —
базовый) означают, что сезонность сдвигает \emph{среднее значение} $y$, а
наклоны при $x_{t1}$, $x_{t2}$, $x_{t3}$ одинаковы во всех кварталах:
\[
  y_t = \beta_1 + \beta_2x_{t1} + \beta_3x_{t2} + \beta_4x_{t3}
        + \beta_5d_{t1} + \beta_6d_{t2} + \beta_7d_{t3} + \varepsilon_t .
\]
Отсутствие сезонности — это $H_0\colon \beta_5 = \beta_6 = \beta_7 = 0$, то есть
$q = 3$ ограничения; в модели без ограничений $k = 7$ и $n - k = 17$. Так как
$\TSS$ у обеих моделей одна и та же,
\[
\begin{aligned}
  \ESS_{\mathrm{R}} - \ESS_{\mathrm{UR}}
  &= \RSS_{\mathrm{UR}} - \RSS_{\mathrm{R}} = 118{,}20 - 110{,}32 = 7{,}88, \\[2pt]
  \ESS_{\mathrm{UR}} &= \TSS - \RSS_{\mathrm{UR}} = 131{,}75 - 118{,}20 = 13{,}55,
\end{aligned}
\]
и по \eqref{eq:restr}
\[
  F = \frac{7{,}88/3}{13{,}55/17} = 3{,}295 > 3{,}197 = F_{0{,}05}(3, 17).
\]
Гипотезу об отсутствии сезонности отвергаем: сезонность есть.

\medskip
\textbf{д)} Это тест Чоу \eqref{eq:chow} с $k = 4$ и $n = 24$:
\[
  \ESS_{\mathrm{R}} = 21{,}43,
  \qquad
  \ESS_1 + \ESS_2 = 12{,}25 + 2{,}32 = 14{,}57,
\]
\[
  F = \frac{(21{,}43 - 14{,}57)/4}{14{,}57/16} = 1{,}883 < 3{,}007 = F_{0{,}05}(4, 16).
\]
Гипотеза об отсутствии структурного изменения не отвергается: данные не дают
оснований утверждать, что между 1-м и 2-м кварталами 1975 г. произошёл перелом.
\end{solution}

% -------------------------------------------------------------------
\problem{4.14}

В таблице представлены совокупный объём внутренних инвестиций $y$ и валовой
внутренний продукт США (млрд долл.) за период с 1939 по 1954 г.

\begin{center}
\begin{tabular}{ccc @{\qquad} ccc}
\toprule
Год & $y$ & $x$ & Год & $y$ & $x$ \\
\midrule
1939 &  9,3 &  90,8 & 1947 & 34,0 & 232,8 \\
1940 & 13,1 & 100,0 & 1948 & 45,9 & 259,1 \\
1941 & 17,9 & 124,9 & 1949 & 35,3 & 258,0 \\
1942 &  9,9 & 158,3 & 1950 & 53,8 & 286,2 \\
1943 &  5,8 & 192,0 & 1951 & 59,2 & 330,2 \\
1944 &  7,2 & 210,5 & 1952 & 52,1 & 347,2 \\
1945 & 10,6 & 212,3 & 1953 & 53,3 & 366,1 \\
1946 & 30,7 & 209,3 & 1954 & 52,7 & 366,3 \\
\bottomrule
\end{tabular}
\end{center}

\noindent Напишите и оцените уравнения, позволяющие ответить на вопрос,
изменилась ли зависимость инвестиций от валового внутреннего продукта во время
войны (1942--1945 гг.) по сравнению с мирным временем.

\begin{excel}
Данные вводятся двумя столбцами, $y$ и $x$. Для модели с фиктивной переменной
добавьте ещё два столбца: \xl{=ЕСЛИ(И(A2>=1942;A2<=1945);1;0)} для $D$ и
произведение \xl{=D2*C2} для $D\cdot x$ (здесь A — год, C — $x$).

\smallskip
Дальше \ui{Данные} $\to$ \ui{Анализ данных} $\to$ \ui{Регрессия}. Чтобы
получить тест Чоу вручную, проведите три регрессии: по всем 16 наблюдениям, по
12 мирным годам и по 4 военным, и возьмите из отчётов строку
\ui{Остаток} столбца \ui{SS} — это $\ESS$.
\end{excel}

\begin{solution}
Военные годы заметно выпадают из общей картины: при растущем ВВП инвестиции в
1942--1945 гг. падают. Проверим это формально.

\medskip
\textbf{1) Регрессия по всем наблюдениям} ($n = 16$, $k = 2$):
\begin{center}
\begin{tabular}{lrrrr}
\toprule
Переменная & Коэффициент & Ст. ошибка & $t$-статистика & $P$-значение \\
\midrule
const & $-14{,}362$ & 7,543 & $-1{,}904$ & 0,0777 \\
$x$   & $\phantom{-}0{,}192$ & 0,030 & $\phantom{-}6{,}368$ & 0,0000 \\
\midrule
\multicolumn{5}{l}{$R^2 = 0{,}7434$;\quad $\ESS_{\mathrm{R}} = 1540{,}33$} \\
\bottomrule
\end{tabular}
\end{center}

\textbf{2) Та же регрессия без военных лет} ($n = 12$):
\begin{center}
\begin{tabular}{lrrrr}
\toprule
Переменная & Коэффициент & Ст. ошибка & $t$-статистика & $P$-значение \\
\midrule
const & $-3{,}335$ & 4,086 & $-0{,}816$ & 0,4334 \\
$x$   & $\phantom{-}0{,}167$ & 0,015 & $10{,}872$ & 0,0000 \\
\midrule
\multicolumn{5}{l}{$R^2 = 0{,}9220$;\quad $\ESS_1 = 260{,}12$} \\
\bottomrule
\end{tabular}
\end{center}

\noindent Во второй регрессии и $t$-статистика наклона, и коэффициент
детерминации заметно выше: без военных лет зависимость гораздо теснее.

\medskip
\textbf{3) Регрессия только по военным годам} ($n = 4$):
\begin{center}
\begin{tabular}{lrrrr}
\toprule
Переменная & Коэффициент & Ст. ошибка & $t$-статистика & $P$-значение \\
\midrule
const & $\phantom{-}11{,}244$ & 12,220 & $\phantom{-}0{,}920$ & 0,4546 \\
$x$   & $-0{,}015$ & \phantom{0}0,063 & $-0{,}236$ & 0,8352 \\
\midrule
\multicolumn{5}{l}{$R^2 = 0{,}0272$;\quad $\ESS_2 = 14{,}87$} \\
\bottomrule
\end{tabular}
\end{center}

\noindent Четыре наблюдения при двух оцениваемых параметрах оставляют
$n - k = 2$ степени свободы, поэтому оценки крайне неточны: наклон меняет знак,
но незначим ($P = 0{,}84$), а $R^2$ практически нулевой. Сама по себе эта
регрессия ничего не доказывает — она нужна, чтобы получить $\ESS_2$ для теста
Чоу. Заодно видно, чем именно различаются периоды: в мирные годы инвестиции
растут вместе с ВВП, в военные связь пропадает.

\medskip
\textbf{4) Тест Чоу.} Разобьём выборку на мирное время (1939--1941, 1946--1954;
$n_1 = 12$) и военное (1942--1945; $n_2 = 4$). Остаточные суммы уже
посчитаны: $\ESS_{\mathrm{R}}$ из п.~1), $\ESS_1$ из п.~2), $\ESS_2$ из п.~3).
По \eqref{eq:chow} при $k = 2$, $n = 16$
\[
  F = \frac{\bigl(\ESS_{\mathrm{R}} - (\ESS_1 + \ESS_2)\bigr)/k}
           {(\ESS_1 + \ESS_2)/(n - 2k)}
    = \frac{\bigl(\ESS_{\mathrm{R}} - (\ESS_1 + \ESS_2)\bigr)/2}
           {(\ESS_1 + \ESS_2)/12},
\]
\[
  F = \frac{\bigl(1540{,}33 - (260{,}12 + 14{,}87)\bigr)/2}
           {(260{,}12 + 14{,}87)/12}
    = \frac{1265{,}34/2}{274{,}99/12}
    = 27{,}608 > 3{,}885 = F_{0{,}05}(2, 12).
\]
Гипотеза о совпадении моделей уверенно отвергается ($P$-значение практически
нулевое): во время войны зависимость инвестиций от ВВП была иной.

\medskip
\textbf{5) То же через фиктивную переменную.} Введём $D_t = 1$ для 1942--1945
гг. и $D_t = 0$ иначе и оценим по всем 16 наблюдениям модель с взаимодействием
\eqref{eq:dummy}:
\[
  y_t = \beta_1 + \beta_2x_t + \beta_3D_t + \beta_4(D_tx_t) + \varepsilon_t .
\]
\begin{center}
\begin{tabular}{lrrrr}
\toprule
Переменная & Коэффициент & Ст. ошибка & $t$-статистика & $P$-значение \\
\midrule
const      & $-3{,}335$ & \phantom{0}3,836 & $-0{,}870$ & 0,4016 \\
$x$        & $\phantom{-}0{,}167$ & \phantom{0}0,014 & $11{,}583$ & 0,0000 \\
$D$        & $14{,}579$ & 21,793 & $\phantom{-}0{,}669$ & 0,5162 \\
$D\cdot x$ & $-0{,}182$ & \phantom{0}0,111 & $-1{,}638$ & 0,1273 \\
\midrule
\multicolumn{5}{l}{$R^2 = 0{,}9542$;\quad $\ESS_{\mathrm{UR}} = 274{,}99$} \\
\bottomrule
\end{tabular}
\end{center}

\noindent Коэффициенты при $x$ и при константе в точности совпали с мирной
регрессией: базовая группа — мирное время, — а суммы
$\widehat{\beta}_1 + \widehat{\beta}_3 = -3{,}335 + 14{,}579 = 11{,}244$ и
$\widehat{\beta}_2 + \widehat{\beta}_4 = 0{,}167 - 0{,}182 = -0{,}015$
воспроизводят военную регрессию из п.~3). Остаточная сумма этой модели равна
$260{,}12 + 14{,}87 = 274{,}99$, то есть сумме остаточных сумм двух раздельных
регрессий, как и сказано в замечании к п.~\ref{sec:chow}. Поэтому $F$-тест
гипотезы $H_0\colon \beta_3 = \beta_4 = 0$ по \eqref{eq:restr} даёт то же
значение $F = 27{,}608$.

\begin{remark}
По отдельности ни $D$ ($P = 0{,}52$), ни $D\cdot x$ ($P = 0{,}13$) не значимы, а
вместе они значимы с $P$-значением порядка $10^{-5}$. Это типичная картина при
коррелированных регрессорах (п.~\ref{sec:vif}): $D$ и $D\cdot x$ связаны почти
функционально, поэтому каждый по отдельности оценивается неточно. Вывод о
структурном сдвиге нужно делать по совместному $F$-тесту, а не по отдельным
$t$-статистикам.
\end{remark}
\end{solution}

% -------------------------------------------------------------------
\ifsolutions\clearpage\fi
\problem{4.15}

В таблице представлены квартальные данные об объёмах продаж и доходах
текстильных корпораций США с первого квартала 1974 г. по третий квартал
1979 г. Введите сезонные фиктивные переменные и с помощью регрессии дохода на
объём продаж исследуйте наличие или отсутствие сезонных колебаний.

\begin{center}
\small
\begin{tabular}{cccc @{\qquad} cccc}
\toprule
Год & Кв. & Продажи & Доход & Год & Кв. & Продажи & Доход \\
\midrule
1974 & I   & 242,0 & 13,5 & 1977 & I   & 311,5 & 15,6 \\
     & II  & 269,4 & 16,3 &      & II  & 338,6 & 19,7 \\
     & III & 272,1 & 15,5 &      & III & 331,7 & 16,7 \\
     & IV  & 277,0 & 13,4 &      & IV  & 346,2 & 18,4 \\
\midrule
1975 & I   & 247,1 & \phantom{0}9,3 & 1978 & I   & 340,2 & 16,0 \\
     & II  & 265,8 & 12,4 &      & II  & 377,5 & 22,1 \\
     & III & 271,0 & 13,2 &      & III & 376,9 & 20,4 \\
     & IV  & 281,3 & 14,2 &      & IV  & 401,8 & 22,6 \\
\midrule
1976 & I   & 284,2 & 14,8 & 1979 & I   & 406,2 & 22,6 \\
     & II  & 307,6 & 18,1 &      & II  & 436,4 & 26,8 \\
     & III & 301,6 & 16,0 &      & III & 437,5 & 24,8 \\
     & IV  & 309,8 & 15,6 &      &     &       &      \\
\bottomrule
\end{tabular}
\end{center}

\begin{excel}
Столбцы \ui{VOL} (продажи) и \ui{INC} (доход) вводятся как есть. Квартальные
фиктивные переменные удобно получить формулой по номеру квартала: если номер
квартала стоит в столбце B, то \xl{=ЕСЛИ(\$B2=1;1;0)} для \ui{QRT1} и так далее.

\smallskip
\textbf{Включайте только три фиктивные переменные из четырёх.} Если добавить все
четыре, их сумма совпадёт со столбцом единиц, и Excel выдаст ошибку или нулевые
коэффициенты — это ловушка фиктивных переменных.
\end{excel}

% Решение с новой страницы: иначе таблицы попадают в конец страницы внутри
% рамки решения и наезжают на колонтитул.
\ifsolutions\clearpage\fi
\begin{solution}
Обозначим через $\mathit{VOL}$ объём продаж, через $\mathit{INC}$ — доход;
$n = 23$ наблюдения.

\medskip
\textbf{1) Регрессия без сезонности.}
\begin{center}
\begin{tabular}{lrrrr}
\toprule
Переменная & Коэффициент & Ст. ошибка & $t$-статистика & $P$-значение \\
\midrule
const          & $-5{,}056$ & 1,682 & $-3{,}006$ & 0,0067 \\
$\mathit{VOL}$ & $\phantom{-}0{,}069$ & 0,005 & $13{,}503$ & 0,0000 \\
\midrule
\multicolumn{5}{l}{$R^2 = 0{,}8967$;\quad $\ESS_{\mathrm{R}} = 42{,}04$;\quad $n - k = 21$} \\
\bottomrule
\end{tabular}
\end{center}

\medskip
\textbf{2) Добавляем сезонные фиктивные переменные.} Вводим $\mathit{QRT}i = 1$,
если наблюдение относится к $i$-му кварталу, и 0 иначе. В регрессию включаем
только $\mathit{QRT}1$, $\mathit{QRT}2$, $\mathit{QRT}3$: четвёртый квартал —
базовый (п.~\ref{sec:dummy}).
\begin{center}
\begin{tabular}{lrrrr}
\toprule
Переменная & Коэффициент & Ст. ошибка & $t$-статистика & $P$-значение \\
\midrule
const             & $-4{,}874$ & 1,568 & $-3{,}109$ & 0,0061 \\
$\mathit{VOL}$    & $\phantom{-}0{,}067$ & 0,005 & $14{,}791$ & 0,0000 \\
$\mathit{QRT}1$   & $-0{,}329$ & 0,749 & $-0{,}440$ & 0,6655 \\
$\mathit{QRT}2$   & $\phantom{-}1{,}767$ & 0,746 & $\phantom{-}2{,}368$ & 0,0293 \\
$\mathit{QRT}3$   & $\phantom{-}0{,}350$ & 0,746 & $\phantom{-}0{,}470$ & 0,6443 \\
\midrule
\multicolumn{5}{l}{$R^2 = 0{,}9331$;\quad $\ESS_{\mathrm{UR}} = 27{,}23$;\quad $n - k = 18$} \\
\bottomrule
\end{tabular}
\end{center}

\noindent По отдельности значим лишь коэффициент при $\mathit{QRT}2$
($P = 0{,}029$). Но вопрос «есть ли сезонность» — это совместная гипотеза
$H_0\colon \gamma_1 = \gamma_2 = \gamma_3 = 0$, то есть $q = 3$ ограничения в
\eqref{eq:restr}. Модель с ограничениями — регрессия из п.~1)
($\ESS_{\mathrm{R}} = 42{,}04$), без ограничений — из п.~2)
($\ESS_{\mathrm{UR}} = 27{,}23$, $n - k = 23 - 5 = 18$):
\[
  F = \frac{\bigl(\ESS_{\mathrm{R}} - \ESS_{\mathrm{UR}}\bigr)/q}
           {\ESS_{\mathrm{UR}}/(n-k)}
    = \frac{\bigl(\ESS_{\mathrm{R}} - \ESS_{\mathrm{UR}}\bigr)/3}
           {\ESS_{\mathrm{UR}}/18},
\]
\[
  F = \frac{(42{,}04 - 27{,}23)/3}{27{,}23/18} = 3{,}262 > 3{,}160 = F_{0{,}05}(3, 18),
  \qquad P = 0{,}046 .
\]
Гипотезу об отсутствии сезонности отвергаем, хотя и на грани.

\medskip
\textbf{3) Какие кварталы отличаются.} Проверим, можно ли убрать
$\mathit{QRT}1$ и $\mathit{QRT}3$, то есть $H_0\colon \gamma_1 = \gamma_3 = 0$
($q = 2$). Модель с ограничением — регрессия на $\mathit{VOL}$ и
$\mathit{QRT}2$, теперь для неё $\ESS_{\mathrm{R}} = 28{,}58$; модель без
ограничений та же, что в п.~2), $\ESS_{\mathrm{UR}} = 27{,}23$. По
\eqref{eq:restr}
\[
  F = \frac{\bigl(\ESS_{\mathrm{R}} - \ESS_{\mathrm{UR}}\bigr)/q}
           {\ESS_{\mathrm{UR}}/(n-k)}
    = \frac{\bigl(\ESS_{\mathrm{R}} - \ESS_{\mathrm{UR}}\bigr)/2}
           {\ESS_{\mathrm{UR}}/18},
\]
\[
  F = \frac{(28{,}58 - 27{,}23)/2}{27{,}23/18} = 0{,}445 < 3{,}555 = F_{0{,}05}(2, 18).
\]
Гипотеза не отвергается: первый и третий кварталы от базового четвёртого не
отличаются. Итоговая модель:
\begin{center}
\begin{tabular}{lrrrr}
\toprule
Переменная & Коэффициент & Ст. ошибка & $t$-статистика & $P$-значение \\
\midrule
const           & $-5{,}096$ & 1,421 & $-3{,}585$ & 0,0018 \\
$\mathit{VOL}$  & $\phantom{-}0{,}068$ & 0,004 & $15{,}610$ & 0,0000 \\
$\mathit{QRT}2$ & $\phantom{-}1{,}750$ & 0,570 & $\phantom{-}3{,}069$ & 0,0061 \\
\midrule
\multicolumn{5}{l}{$R^2 = 0{,}9298$;\quad $R^2_{\text{adj}} = 0{,}9228$} \\
\bottomrule
\end{tabular}
\end{center}

\noindent\textbf{Вывод.} Сезонность есть, но сводится к одному кварталу: при
одном и том же объёме продаж доход во втором квартале в среднем выше на
1,75 млн долл., чем в остальных. Заметим, что $R^2$ при переходе от модели с
тремя фиктивными переменными к модели с одной снизился ($0{,}9331$ против
$0{,}9298$), а скорректированный коэффициент \eqref{eq:r2adj} вырос
($0{,}9182$ против $0{,}9228$): по критерию \eqref{eq:r2adj} экономная модель
лучше.
\end{solution}
